% Options for packages loaded elsewhere
% Options for packages loaded elsewhere
\PassOptionsToPackage{unicode}{hyperref}
\PassOptionsToPackage{hyphens}{url}
\PassOptionsToPackage{dvipsnames,svgnames,x11names}{xcolor}
%
\documentclass[
  letterpaper,
]{report}
\usepackage{xcolor}
\usepackage{amsmath,amssymb}
\setcounter{secnumdepth}{5}
\usepackage{iftex}
\ifPDFTeX
  \usepackage[T1]{fontenc}
  \usepackage[utf8]{inputenc}
  \usepackage{textcomp} % provide euro and other symbols
\else % if luatex or xetex
  \usepackage{unicode-math} % this also loads fontspec
  \defaultfontfeatures{Scale=MatchLowercase}
  \defaultfontfeatures[\rmfamily]{Ligatures=TeX,Scale=1}
\fi
\usepackage{lmodern}
\ifPDFTeX\else
  % xetex/luatex font selection
\fi
% Use upquote if available, for straight quotes in verbatim environments
\IfFileExists{upquote.sty}{\usepackage{upquote}}{}
\IfFileExists{microtype.sty}{% use microtype if available
  \usepackage[]{microtype}
  \UseMicrotypeSet[protrusion]{basicmath} % disable protrusion for tt fonts
}{}
\makeatletter
\@ifundefined{KOMAClassName}{% if non-KOMA class
  \IfFileExists{parskip.sty}{%
    \usepackage{parskip}
  }{% else
    \setlength{\parindent}{0pt}
    \setlength{\parskip}{6pt plus 2pt minus 1pt}}
}{% if KOMA class
  \KOMAoptions{parskip=half}}
\makeatother
% Make \paragraph and \subparagraph free-standing
\makeatletter
\ifx\paragraph\undefined\else
  \let\oldparagraph\paragraph
  \renewcommand{\paragraph}{
    \@ifstar
      \xxxParagraphStar
      \xxxParagraphNoStar
  }
  \newcommand{\xxxParagraphStar}[1]{\oldparagraph*{#1}\mbox{}}
  \newcommand{\xxxParagraphNoStar}[1]{\oldparagraph{#1}\mbox{}}
\fi
\ifx\subparagraph\undefined\else
  \let\oldsubparagraph\subparagraph
  \renewcommand{\subparagraph}{
    \@ifstar
      \xxxSubParagraphStar
      \xxxSubParagraphNoStar
  }
  \newcommand{\xxxSubParagraphStar}[1]{\oldsubparagraph*{#1}\mbox{}}
  \newcommand{\xxxSubParagraphNoStar}[1]{\oldsubparagraph{#1}\mbox{}}
\fi
\makeatother

\usepackage{color}
\usepackage{fancyvrb}
\newcommand{\VerbBar}{|}
\newcommand{\VERB}{\Verb[commandchars=\\\{\}]}
\DefineVerbatimEnvironment{Highlighting}{Verbatim}{commandchars=\\\{\}}
% Add ',fontsize=\small' for more characters per line
\usepackage{framed}
\definecolor{shadecolor}{RGB}{241,243,245}
\newenvironment{Shaded}{\begin{snugshade}}{\end{snugshade}}
\newcommand{\AlertTok}[1]{\textcolor[rgb]{0.68,0.00,0.00}{#1}}
\newcommand{\AnnotationTok}[1]{\textcolor[rgb]{0.37,0.37,0.37}{#1}}
\newcommand{\AttributeTok}[1]{\textcolor[rgb]{0.40,0.45,0.13}{#1}}
\newcommand{\BaseNTok}[1]{\textcolor[rgb]{0.68,0.00,0.00}{#1}}
\newcommand{\BuiltInTok}[1]{\textcolor[rgb]{0.00,0.23,0.31}{#1}}
\newcommand{\CharTok}[1]{\textcolor[rgb]{0.13,0.47,0.30}{#1}}
\newcommand{\CommentTok}[1]{\textcolor[rgb]{0.37,0.37,0.37}{#1}}
\newcommand{\CommentVarTok}[1]{\textcolor[rgb]{0.37,0.37,0.37}{\textit{#1}}}
\newcommand{\ConstantTok}[1]{\textcolor[rgb]{0.56,0.35,0.01}{#1}}
\newcommand{\ControlFlowTok}[1]{\textcolor[rgb]{0.00,0.23,0.31}{\textbf{#1}}}
\newcommand{\DataTypeTok}[1]{\textcolor[rgb]{0.68,0.00,0.00}{#1}}
\newcommand{\DecValTok}[1]{\textcolor[rgb]{0.68,0.00,0.00}{#1}}
\newcommand{\DocumentationTok}[1]{\textcolor[rgb]{0.37,0.37,0.37}{\textit{#1}}}
\newcommand{\ErrorTok}[1]{\textcolor[rgb]{0.68,0.00,0.00}{#1}}
\newcommand{\ExtensionTok}[1]{\textcolor[rgb]{0.00,0.23,0.31}{#1}}
\newcommand{\FloatTok}[1]{\textcolor[rgb]{0.68,0.00,0.00}{#1}}
\newcommand{\FunctionTok}[1]{\textcolor[rgb]{0.28,0.35,0.67}{#1}}
\newcommand{\ImportTok}[1]{\textcolor[rgb]{0.00,0.46,0.62}{#1}}
\newcommand{\InformationTok}[1]{\textcolor[rgb]{0.37,0.37,0.37}{#1}}
\newcommand{\KeywordTok}[1]{\textcolor[rgb]{0.00,0.23,0.31}{\textbf{#1}}}
\newcommand{\NormalTok}[1]{\textcolor[rgb]{0.00,0.23,0.31}{#1}}
\newcommand{\OperatorTok}[1]{\textcolor[rgb]{0.37,0.37,0.37}{#1}}
\newcommand{\OtherTok}[1]{\textcolor[rgb]{0.00,0.23,0.31}{#1}}
\newcommand{\PreprocessorTok}[1]{\textcolor[rgb]{0.68,0.00,0.00}{#1}}
\newcommand{\RegionMarkerTok}[1]{\textcolor[rgb]{0.00,0.23,0.31}{#1}}
\newcommand{\SpecialCharTok}[1]{\textcolor[rgb]{0.37,0.37,0.37}{#1}}
\newcommand{\SpecialStringTok}[1]{\textcolor[rgb]{0.13,0.47,0.30}{#1}}
\newcommand{\StringTok}[1]{\textcolor[rgb]{0.13,0.47,0.30}{#1}}
\newcommand{\VariableTok}[1]{\textcolor[rgb]{0.07,0.07,0.07}{#1}}
\newcommand{\VerbatimStringTok}[1]{\textcolor[rgb]{0.13,0.47,0.30}{#1}}
\newcommand{\WarningTok}[1]{\textcolor[rgb]{0.37,0.37,0.37}{\textit{#1}}}

\usepackage{longtable,booktabs,array}
\newcounter{none} % for unnumbered tables
\usepackage{calc} % for calculating minipage widths
% Correct order of tables after \paragraph or \subparagraph
\usepackage{etoolbox}
\makeatletter
\patchcmd\longtable{\par}{\if@noskipsec\mbox{}\fi\par}{}{}
\makeatother
% Allow footnotes in longtable head/foot
\IfFileExists{footnotehyper.sty}{\usepackage{footnotehyper}}{\usepackage{footnote}}
\makesavenoteenv{longtable}
\usepackage{graphicx}
\makeatletter
\newsavebox\pandoc@box
\newcommand*\pandocbounded[1]{% scales image to fit in text height/width
  \sbox\pandoc@box{#1}%
  \Gscale@div\@tempa{\textheight}{\dimexpr\ht\pandoc@box+\dp\pandoc@box\relax}%
  \Gscale@div\@tempb{\linewidth}{\wd\pandoc@box}%
  \ifdim\@tempb\p@<\@tempa\p@\let\@tempa\@tempb\fi% select the smaller of both
  \ifdim\@tempa\p@<\p@\scalebox{\@tempa}{\usebox\pandoc@box}%
  \else\usebox{\pandoc@box}%
  \fi%
}
% Set default figure placement to htbp
\def\fps@figure{htbp}
\makeatother





\setlength{\emergencystretch}{3em} % prevent overfull lines

\providecommand{\tightlist}{%
  \setlength{\itemsep}{0pt}\setlength{\parskip}{0pt}}



 


% ODK Collect - treeCensus
% \documentclass[12pt, a4paper]{article} %padrao letterpaper, 10pt
\usepackage[utf8]{inputenc}
\usepackage[brazil]{babel}
\usepackage[defaultfam,tabular,lining]{montserrat} %% Option 'defaultfam'
%% only if the base font of the document is to be sans serif
\usepackage[T1]{fontenc}
\renewcommand*\oldstylenums[1]{{\fontfamily{Montserrat-TOsF}\selectfont #1}}
%% font Nimbus Roman ~ Times New Roman
\usepackage{mathptmx}
%\usepackage[latin1]{inputenc}
% \usepackage{dejavu}
% \renewcommand*\familydefault{\ttdefault}
%\usepackage[T1]{fontenc}
% \renewcommand*\familydefault{\sfdefault}
%\usepackage{amsfonts,amssymb,graphicx,enumerate}
% amssymb package, useful for mathematical symbols
%\usepackage[centertags]{amsmath}
\usepackage{caption}
\usepackage{subcaption}
\usepackage[sectionbib]{natbib}
% \bibpunct{(}{)}{;}{a}{}{}
\bibpunct{(}{)}{;}{a}{,}{,}
\usepackage[nottoc,numbib]{tocbibind}

% \usepackage{etoolbox}
% \patchcmd{\thebibliography}{\chapter*}{\section*}{}{}
\usepackage{enumitem}
\setlist{noitemsep}
\usepackage[usenames,dvipsnames, table, xcdraw]{xcolor}
%\definecolor{light-gray}{gray}{0.9}
% Configuracoes de pagina
\usepackage[lmargin=2.5cm,rmargin=2cm,tmargin=2.5cm,bmargin=2.5cm]{geometry}
% Layout da pagina
\usepackage{hyperref}
\hypersetup{pdfnewwindow=true,
    pdfpagelayout=SinglePage, % ou TwoPageLeft
    colorlinks=true,
    citecolor=Violet,
    linkcolor=blue,
    filecolor=ForestGreen,
    urlcolor=ForestGreen,
    pdftitle={Manual ODK Collect PPSP},
    pdfauthor={Alexandre Adalardo de Oliveira}}
%**********************************************
\usepackage{tikz}
\usepackage{float}
\usepackage{pdfpages}
%********* boxes
% Definindo novos comandos
\newcommand{\HRule}{\rule{\linewidth}{0.5mm}}
% \providecommand{\sin}{} \renewcommand{\sin}{\hspace{2pt}\textrm{sen}}
% \newcommand{\R}{\mathbb{R}} %simbolos de numeros 
% \title(Relatório 
% \author{R\'egis da Silva}
% \date{2011}    %para ocultar a data digite: \date{ }
\linespread{1.3}
\usepackage{listings}
\usepackage{graphicx}
\graphicspath{{../media/}}
\usepackage{sidecap}
\usepackage{subcaption}
\date{\today}
\title{}
\author{}
%%%%%%%%%%%%%%%%%%%%%%%%%%%%%%%%%%%%%%%%%%%%%%%%%%%%%
\makeatletter
\@ifpackageloaded{tcolorbox}{}{\usepackage[skins,breakable]{tcolorbox}}
\@ifpackageloaded{fontawesome5}{}{\usepackage{fontawesome5}}
\definecolor{quarto-callout-color}{HTML}{909090}
\definecolor{quarto-callout-note-color}{HTML}{0758E5}
\definecolor{quarto-callout-important-color}{HTML}{CC1914}
\definecolor{quarto-callout-warning-color}{HTML}{EB9113}
\definecolor{quarto-callout-tip-color}{HTML}{00A047}
\definecolor{quarto-callout-caution-color}{HTML}{FC5300}
\definecolor{quarto-callout-color-frame}{HTML}{acacac}
\definecolor{quarto-callout-note-color-frame}{HTML}{4582ec}
\definecolor{quarto-callout-important-color-frame}{HTML}{d9534f}
\definecolor{quarto-callout-warning-color-frame}{HTML}{f0ad4e}
\definecolor{quarto-callout-tip-color-frame}{HTML}{02b875}
\definecolor{quarto-callout-caution-color-frame}{HTML}{fd7e14}
\makeatother
\makeatletter
\@ifpackageloaded{bookmark}{}{\usepackage{bookmark}}
\makeatother
\makeatletter
\@ifpackageloaded{caption}{}{\usepackage{caption}}
\AtBeginDocument{%
\ifdefined\contentsname
  \renewcommand*\contentsname{Table of contents}
\else
  \newcommand\contentsname{Table of contents}
\fi
\ifdefined\listfigurename
  \renewcommand*\listfigurename{List of Figures}
\else
  \newcommand\listfigurename{List of Figures}
\fi
\ifdefined\listtablename
  \renewcommand*\listtablename{List of Tables}
\else
  \newcommand\listtablename{List of Tables}
\fi
\ifdefined\figurename
  \renewcommand*\figurename{Figure}
\else
  \newcommand\figurename{Figure}
\fi
\ifdefined\tablename
  \renewcommand*\tablename{Table}
\else
  \newcommand\tablename{Table}
\fi
}
\@ifpackageloaded{float}{}{\usepackage{float}}
\floatstyle{ruled}
\@ifundefined{c@chapter}{\newfloat{codelisting}{h}{lop}}{\newfloat{codelisting}{h}{lop}[chapter]}
\floatname{codelisting}{Listing}
\newcommand*\listoflistings{\listof{codelisting}{List of Listings}}
\makeatother
\makeatletter
\makeatother
\makeatletter
\@ifpackageloaded{caption}{}{\usepackage{caption}}
\@ifpackageloaded{subcaption}{}{\usepackage{subcaption}}
\makeatother
\usepackage{bookmark}
\IfFileExists{xurl.sty}{\usepackage{xurl}}{} % add URL line breaks if available
\urlstyle{same}
% fallback for those not using the hyperref driver hyperxmp:
\makeatletter
\@ifundefined{xmpquote}{\newcommand{\xmpquote}[1]{#1}}{}
\makeatother
\hypersetup{
  pdftitle={ODK Collect treeCensus Manual},
  pdfauthor={Alexandre Adalardo de Oliveira},
  colorlinks=true,
  linkcolor={blue},
  filecolor={Maroon},
  citecolor={Blue},
  urlcolor={Blue},
  pdfcreator={LaTeX via pandoc}}


\title{ODK Collect treeCensus Manual}
\author{Alexandre Adalardo de Oliveira}
\date{2026-10-01}
\begin{document}
\maketitle


\begin{titlepage}
\begin{center}

% Upper part of the page. The '~' is needed because \\
% only works if a paragraph has started.

\textsc{\large Manual odkTreeCensus}\\
\vspace{3cm}
% Title
\HRule \\
{\Large \bfseries Pacote do R para coleta de dados em campo por dispositivo móvel utilizando o ODK Collect\\}

\HRule 

\vspace{5cm}
\begin{center}
\begin{tikzpicture}
  \node[fill=gray!10!white,rounded corners=5pt,text width=1\textwidth] {
    \textsc{Contexto}\par

    \vspace{.5cm}
O procedimento de coleta de dados em campo tem fontes de incertezas e erros envolvidas na de aquisição da informação. Entre as fonte de erro mais importantes nos inventários florestais estão: erros na anotação do dado e transcrição da informação em papel para o formato digital. Esses erros são custosos para serem resolvidos, muitas vezes demandando o retorno à campo, o que nem sempre é possível. Evitar fontes de erros no procedimento científico é essencial para não comprometer as conclusões dos trabalhos científicos. O pacote \textbf{odkTreeCensus} tem como objetivos minimizar a entrada de dados com discrepâncias durante a aquisição da informação e eliminar a fase de transcrição de dados anotados em papel. O pacote auxilia na construção e gerenciamento de formulários de entrada de dados para o aplicativo ODK-Collect, criando os arquivos necessários e gerenciando a entrada e saída de dados.    

};
\end{tikzpicture}
\end{center}


\vspace{1cm}



{
\begin{flushright} 
Autor:
Alexandre Adalardo de \textsc{Oliveira}\\
{Departamento de Ecologia - IBUSP}
%\end{minipage}
%\begin{minipage}{0.5\textwidth}
\end{flushright}
}
% \end{minipage}
%\vspace{2cm}
% Bottom of the page
%\HRule\\[0.1cm]
\includegraphics[width=0.3\textwidth]{media/LogoLabtrop3.png}



%\pagebreak

%\includegraphics[width=0.75\textwidth]{Arquivos/restingaPq.png}

% % Author and supervisor
% %\begin{minipage}{0.7\textwidth}
% {
% \begin{flushright} 
% Proponente:
% Alexandre Adalardo de \textsc{Oliveira}\\
% {Universidade de São Paulo, }{Departamento de Ecologia - IB}
% %\end{minipage}
% %\begin{minipage}{0.5\textwidth}
% \end{flushright}
% }
% % \end{minipage}
% %\vspace{2cm}
% % Bottom of the page
% %\HRule\\[0.1cm]
% \includegraphics[width=0.3\textwidth]{LogoLabtrop3.png}


\end{center}
\end{titlepage}

\renewcommand*\contentsname{Table of contents}
{
\setcounter{tocdepth}{2}
\tableofcontents
}

\bookmarksetup{startatroot}

\chapter*{Apresentação}\label{apresentauxe7uxe3o}
\addcontentsline{toc}{chapter}{Apresentação}

\markboth{Apresentação}{Apresentação}

Este é o manual de utilização do pacote R \textbf{treeCensus} para
coleta de dados via ODK Collect.

\bookmarksetup{startatroot}

\chapter{Instalando pacote
odkTreeCensus}\label{instalando-pacote-odktreecensus}

 
\begin{titlepage}
\begin{center}

% Upper part of the page. The '~' is needed because \\
% only works if a paragraph has started.

\textsc{\large Manual odkTreeCensus}\\
\vspace{3cm}
% Title
\HRule \\
{\Large \bfseries Pacote do R para coleta de dados em campo por dispositivo móvel utilizando o ODK Collect\\}

\HRule 

\vspace{5cm}
\begin{center}
\begin{tikzpicture}
  \node[fill=gray!10!white,rounded corners=5pt,text width=1\textwidth] {
    \textsc{Contexto}\par

    \vspace{.5cm}
O procedimento de coleta de dados em campo tem fontes de incertezas e erros envolvidas na de aquisição da informação. Entre as fonte de erro mais importantes nos inventários florestais estão: erros na anotação do dado e transcrição da informação em papel para o formato digital. Esses erros são custosos para serem resolvidos, muitas vezes demandando o retorno à campo, o que nem sempre é possível. Evitar fontes de erros no procedimento científico é essencial para não comprometer as conclusões dos trabalhos científicos. O pacote \textbf{odkTreeCensus} tem como objetivos minimizar a entrada de dados com discrepâncias durante a aquisição da informação e eliminar a fase de transcrição de dados anotados em papel. O pacote auxilia na construção e gerenciamento de formulários de entrada de dados para o aplicativo ODK-Collect, criando os arquivos necessários e gerenciando a entrada e saída de dados.    

};
\end{tikzpicture}
\end{center}


\vspace{1cm}



{
\begin{flushright} 
Autor:
Alexandre Adalardo de \textsc{Oliveira}\\
{Departamento de Ecologia - IBUSP}
%\end{minipage}
%\begin{minipage}{0.5\textwidth}
\end{flushright}
}
% \end{minipage}
%\vspace{2cm}
% Bottom of the page
%\HRule\\[0.1cm]
\includegraphics[width=0.3\textwidth]{media/LogoLabtrop3.png}



%\pagebreak

%\includegraphics[width=0.75\textwidth]{Arquivos/restingaPq.png}

% % Author and supervisor
% %\begin{minipage}{0.7\textwidth}
% {
% \begin{flushright} 
% Proponente:
% Alexandre Adalardo de \textsc{Oliveira}\\
% {Universidade de São Paulo, }{Departamento de Ecologia - IB}
% %\end{minipage}
% %\begin{minipage}{0.5\textwidth}
% \end{flushright}
% }
% % \end{minipage}
% %\vspace{2cm}
% % Bottom of the page
% %\HRule\\[0.1cm]
% \includegraphics[width=0.3\textwidth]{LogoLabtrop3.png}


\end{center}
\end{titlepage}


\begin{Shaded}
\begin{Highlighting}[]
\DocumentationTok{\#\# verifica se tem o pacote remotes e caso nao tenha instala}
\ControlFlowTok{if}\NormalTok{(}\SpecialCharTok{!}\StringTok{"remotes"} \SpecialCharTok{\%in\%} \FunctionTok{rownames}\NormalTok{(}\FunctionTok{installed.packages}\NormalTok{()))}
\NormalTok{\{}
    \FunctionTok{install.package}\NormalTok{(}\StringTok{"remote"}\NormalTok{)}
\NormalTok{\}}
\DocumentationTok{\#\# Verifica se o odkTreeCensus está carregado na sessão}
\ControlFlowTok{if}\NormalTok{(}\StringTok{"package:odkTreeCensus"} \SpecialCharTok{\%in\%} \FunctionTok{search}\NormalTok{())}
\NormalTok{\{}
    \FunctionTok{detach}\NormalTok{(}\StringTok{"package:odkTreeCensus"}\NormalTok{, }\AttributeTok{unload =} \ConstantTok{TRUE}\NormalTok{)}
\NormalTok{\}}
\DocumentationTok{\#\# verifica sem tem o odkTreeCensus e desinstala}
\ControlFlowTok{if}\NormalTok{(}\StringTok{"odkTreeCensus"} \SpecialCharTok{\%in\%} \FunctionTok{rownames}\NormalTok{(}\FunctionTok{installed.packages}\NormalTok{()))}
\NormalTok{\{}
    \FunctionTok{remove.packages}\NormalTok{(}\StringTok{"odkTreeCensus"}\NormalTok{)}
\NormalTok{\}}
\DocumentationTok{\#\# instala a versao recente do odkTreeCensus}

\NormalTok{remotes}\SpecialCharTok{::}\FunctionTok{install\_github}\NormalTok{(}\StringTok{"adalardo/odkTreeCensus"}\NormalTok{)}
\FunctionTok{library}\NormalTok{(}\StringTok{"odkTreeCensus"}\NormalTok{)}
\end{Highlighting}
\end{Shaded}

\section{Criar arquivos de dados
acessórios}\label{criar-arquivos-de-dados-acessuxf3rios}

Todos os arquivos de dados que acompanham o formulário no ODK Collect
devem ser arquivos em texto puro no formato csv, separado por ``,'' e
decimal ``.''. Eles são criados em um diretorio local no computador para
depois serem transferidos para o \texttt{ODK\ Collect}.

\begin{tcolorbox}[enhanced jigsaw, arc=.35mm, bottomrule=.15mm, bottomtitle=1mm, breakable, colback=white, colbacktitle=quarto-callout-tip-color!10!white, colframe=quarto-callout-tip-color-frame, coltitle=black, left=2mm, leftrule=.75mm, opacityback=0, opacitybacktitle=0.6, rightrule=.15mm, title=\textcolor{quarto-callout-tip-color}{\faLightbulb}\hspace{0.5em}{Importante}, titlerule=0mm, toprule=.15mm, toptitle=1mm]

ude o caminho do objeto \texttt{dirODK} para um diretório local do seu
computador onde pretende salvar os arquivos de media do
\texttt{OKD\ Collect}. Este vai se o diretório de onde os arquivos sẽrão
transferidos do computador para o dispositivo móvel com o
\texttt{ODK\ Collect}.

\end{tcolorbox}

\begin{Shaded}
\begin{Highlighting}[]
\NormalTok{dirODK }\OtherTok{\textless{}{-}} \StringTok{"/home/aao/Ale2026/AleProjetos/Rpackages/odkData"}
\end{Highlighting}
\end{Shaded}

\subsection{Dados das árvores do último censo e códigos de
subparcelas}\label{dados-das-uxe1rvores-do-uxfaltimo-censo-e-cuxf3digos-de-subparcelas}

O \texttt{lastcensus.csv} é um arquivo de dados do censo anterior. Veja
o \texttt{help(lastCensusODK)}. Os dados \texttt{peicCensus2019} são
distribuídos junto ao pacote em um arquivo \texttt{.rda}.

\begin{Shaded}
\begin{Highlighting}[]
\CommentTok{\#help(lastCensusODK)}

\DocumentationTok{\#\# peicCensus2019 \textless{}{-} read.table(file = "/home/aao/Ale2026/AleProjetos/PPEIC/dados/dadosCenso2019/censoPeic2019.csv", sep = ",", header = TRUE)}
\DocumentationTok{\#\# str(peic2019)}
\DocumentationTok{\#\# save(peicCensus2019, file = "data/lastCensusPEIC.rda", compress = "bzip2")}
\DocumentationTok{\#\# dir("data")}
\FunctionTok{data}\NormalTok{(}\AttributeTok{package =} \StringTok{"odkTreeCensus"}\NormalTok{)}
\FunctionTok{data}\NormalTok{(}\StringTok{"lastCensusPEIC"}\NormalTok{, }\AttributeTok{package =} \StringTok{"odkTreeCensus"}\NormalTok{)}
\FunctionTok{str}\NormalTok{(peicCensus2019)}
\NormalTok{lastCensus }\OtherTok{\textless{}{-}} \FunctionTok{lastCensusODK}\NormalTok{(}\AttributeTok{data =}\NormalTok{ peicCensus2019 , }\AttributeTok{splitX =} \DecValTok{5}\NormalTok{, }\AttributeTok{splitY =} \DecValTok{5}\NormalTok{, }\AttributeTok{maxX =} \DecValTok{20}\NormalTok{, }\AttributeTok{maxY =} \DecValTok{20}\NormalTok{, }\AttributeTok{subPlotCodeX =}\NormalTok{ xCode, }\AttributeTok{subPlotCodeY =}\NormalTok{ yCode, }\AttributeTok{dirData =}\NormalTok{ dirODK, }\AttributeTok{save =} \ConstantTok{TRUE}\NormalTok{)}
\FunctionTok{str}\NormalTok{(lastCensus)}
\end{Highlighting}
\end{Shaded}

\section{Códigos das subparcelas}\label{cuxf3digos-das-subparcelas}

Os arquivos \texttt{xcode.csv} e \texttt{ycode.csv} armazenam os códigos
das subparcelas do plot e tem duas colunas \texttt{label} e
\texttt{name}, que identificam as subparcelas no eixo x e y,
respectivamente. As colunas podem ter a mesma informação, mas ambas são
necessárias.

\begin{Shaded}
\begin{Highlighting}[]
\NormalTok{xCode }\OtherTok{\textless{}{-}}\NormalTok{  LETTERS[}\DecValTok{1}\SpecialCharTok{:}\DecValTok{16}\NormalTok{]}
\NormalTok{yCode }\OtherTok{\textless{}{-}} \FunctionTok{sprintf}\NormalTok{(}\StringTok{"\%02d"}\NormalTok{, }\DecValTok{0}\SpecialCharTok{:}\DecValTok{15}\NormalTok{)}
\FunctionTok{write.table}\NormalTok{(}\FunctionTok{data.frame}\NormalTok{(}\AttributeTok{name =}\NormalTok{ xCode, }\AttributeTok{label =}\NormalTok{ xCode), }\AttributeTok{file =} \FunctionTok{file.path}\NormalTok{(dirODK,}\StringTok{"xcode.csv"}\NormalTok{),  }\AttributeTok{sep =} \StringTok{\textquotesingle{},\textquotesingle{}}\NormalTok{, }\AttributeTok{row.names =} \ConstantTok{FALSE}\NormalTok{)}
\FunctionTok{write.table}\NormalTok{(}\FunctionTok{data.frame}\NormalTok{(}\AttributeTok{name =}\NormalTok{ yCode, }\AttributeTok{label =}\NormalTok{ yCode), }\AttributeTok{file =} \FunctionTok{file.path}\NormalTok{(dirODK,}\StringTok{"ycode.csv"}\NormalTok{),  }\AttributeTok{sep =} \StringTok{\textquotesingle{},\textquotesingle{}}\NormalTok{, }\AttributeTok{row.names =} \ConstantTok{FALSE}\NormalTok{)}
\end{Highlighting}
\end{Shaded}

\subsection{Equipe}\label{equipe}

O arquivo \texttt{team.csv} contém duas colunas: \texttt{name} e
\texttt{label} nome e apelido dos membros da equipe de campo.

\begin{Shaded}
\begin{Highlighting}[]
\NormalTok{lb }\OtherTok{\textless{}{-}} \FunctionTok{c}\NormalTok{(}\StringTok{"Alexandre Adalardo"}\NormalTok{, }\StringTok{"Marcel Zimmermann"}\NormalTok{, }\StringTok{"Luisa Truffi"}\NormalTok{)}
\NormalTok{nm }\OtherTok{\textless{}{-}} \FunctionTok{c}\NormalTok{(}\StringTok{"AAO"}\NormalTok{, }\StringTok{"GIL"}\NormalTok{, }\StringTok{"LU"}\NormalTok{)}

\FunctionTok{write.table}\NormalTok{(}\FunctionTok{data.frame}\NormalTok{(}\AttributeTok{name =}\NormalTok{ nm, }\AttributeTok{label =}\NormalTok{ lb), }\AttributeTok{file =} \FunctionTok{file.path}\NormalTok{(dirODK,}\StringTok{"team.csv"}\NormalTok{),  }\AttributeTok{sep =} \StringTok{\textquotesingle{},\textquotesingle{}}\NormalTok{, }\AttributeTok{row.names =} \ConstantTok{FALSE}\NormalTok{)}
\end{Highlighting}
\end{Shaded}

\subsection{Lista de Espécies}\label{lista-de-espuxe9cies}

A lista de espécies contempla todos os nomes, válidos, morfotipos e
\texttt{indet} para todos os níveis taxonômicos a partir de familia. O
arquivo deve-se chamar \texttt{splist.csv}. A função \texttt{spList()}
cria no formato adequado à partir de um arquivo \texttt{fam} e
\texttt{species}. Gera \texttt{gen} se não houver.

\begin{Shaded}
\begin{Highlighting}[]
\FunctionTok{str}\NormalTok{(peic19)}
\NormalTok{splist }\OtherTok{\textless{}{-}} \FunctionTok{spList}\NormalTok{(}\AttributeTok{data =}\NormalTok{ odkdata, }\AttributeTok{dirData =}\NormalTok{ dirODK, }\AttributeTok{saveFile =} \ConstantTok{TRUE}\NormalTok{)}
\FunctionTok{str}\NormalTok{(splist)}
\end{Highlighting}
\end{Shaded}

\subsection{Dados de tags por
subquadrat}\label{dados-de-tags-por-subquadrat}

Essa informação é importante para a lista de tags não encontradas, o
formato é de uma coluna \texttt{subq} como o código do subquadrat e
\texttt{tags} contendo todas as tags do subquadrat separadas por espaço
no mesmo campo.

\begin{Shaded}
\begin{Highlighting}[]
\FunctionTok{subqTags}\NormalTok{(}\AttributeTok{data =}\NormalTok{ odkdata, }\AttributeTok{dirData =}\NormalTok{ dirODK, }\AttributeTok{saveFile =} \ConstantTok{TRUE}\NormalTok{)}
\end{Highlighting}
\end{Shaded}

\section{Mapas e Grids svg}\label{mapas-e-grids-svg}

Os mapas e grids do formulário foram feitos em \texttt{svg} para
permitir a seleção de elementos na imagem. Para isso usamos os pacotes
\texttt{grid} e \texttt{gridSVG}.

\subsection{Os rótulos para os svg}\label{os-ruxf3tulos-para-os-svg}

Para que os elementos do \texttt{svg} sejam selecionados é preciso que
os rótulos sejam opções na planilha \texttt{choices} do formulário
\texttt{xlsx}. Quando há muitas opções é preferível incorporá-las em um
arquivo \texttt{csv}.

Como os mapas incorporam uma borda além do \texttt{subquadrat} é preciso
fazer um objeto de dados de mapas, completo. Esse arquivo é uma lista
grande quando é feito para a parcela toda, mas pode ser feito
parcialmente para subparcelas.

\begin{tcolorbox}[enhanced jigsaw, arc=.35mm, bottomrule=.15mm, bottomtitle=1mm, breakable, colback=white, colbacktitle=quarto-callout-note-color!10!white, colframe=quarto-callout-note-color-frame, coltitle=black, left=2mm, leftrule=.75mm, opacityback=0, opacitybacktitle=0.6, rightrule=.15mm, title=\textcolor{quarto-callout-note-color}{\faInfo}\hspace{0.5em}{Árvores vivas}, titlerule=0mm, toprule=.15mm, toptitle=1mm]

O arquivo está sendo contruído apenas com as árvores vivas no censo
anterior, apesar de ser desejável ter árvores morta e desaparecidas,
isso provavelmente,deixaria o mapa muito poluído.

\end{tcolorbox}

\begin{Shaded}
\begin{Highlighting}[]
\NormalTok{subplot }\OtherTok{\textless{}{-}} \FunctionTok{sort}\NormalTok{(}\FunctionTok{unique}\NormalTok{(odkdata}\SpecialCharTok{$}\NormalTok{quad))}
\NormalTok{mapdata }\OtherTok{\textless{}{-}} \FunctionTok{mapData}\NormalTok{(odkdata[odkdata}\SpecialCharTok{$}\NormalTok{status }\SpecialCharTok{==} \StringTok{"A"}\NormalTok{, ], }\AttributeTok{subplotCodes =}\NormalTok{ subplot, }\AttributeTok{splitX =} \DecValTok{5}\NormalTok{, }\AttributeTok{splitY =} \DecValTok{5}\NormalTok{, }\AttributeTok{maxX =} \DecValTok{20}\NormalTok{, }\AttributeTok{maxY =} \DecValTok{20}\NormalTok{, }\AttributeTok{buffer =} \DecValTok{2}\NormalTok{)}
\NormalTok{sk }\OtherTok{\textless{}{-}} \FunctionTok{subqKey}\NormalTok{(}\AttributeTok{dataMap =}\NormalTok{ mapdata, }\AttributeTok{dirData =}\NormalTok{ dirODK, }\AttributeTok{saveFile =} \ConstantTok{TRUE}\NormalTok{)}
\FunctionTok{str}\NormalTok{(sk)}
\NormalTok{sk[sk}\SpecialCharTok{$}\NormalTok{subq }\SpecialCharTok{==} \StringTok{"A00\_0x0"}\NormalTok{, ]}
\end{Highlighting}
\end{Shaded}

\subsection{gridid.csv}\label{gridid.csv}

\begin{Shaded}
\begin{Highlighting}[]
\NormalTok{subqSize }\OtherTok{=} \DecValTok{5}
\NormalTok{gridSize }\OtherTok{=}\NormalTok{ .}\DecValTok{2}
\NormalTok{xseq }\OtherTok{\textless{}{-}}  \FunctionTok{rep}\NormalTok{(}\FunctionTok{seq}\NormalTok{(}\FloatTok{0.0}\NormalTok{, subqSize }\SpecialCharTok{{-}}\NormalTok{ gridSize, }\AttributeTok{by =}\NormalTok{ gridSize), }\AttributeTok{each =}\NormalTok{ subqSize}\SpecialCharTok{/}\NormalTok{gridSize)}
\NormalTok{yseq }\OtherTok{\textless{}{-}}  \FunctionTok{rep}\NormalTok{(}\FunctionTok{seq}\NormalTok{(}\FloatTok{0.0}\NormalTok{, subqSize }\SpecialCharTok{{-}}\NormalTok{ gridSize, }\AttributeTok{by =}\NormalTok{ gridSize), subqSize}\SpecialCharTok{/}\NormalTok{gridSize)}

\NormalTok{xyid  }\OtherTok{\textless{}{-}}  \FunctionTok{paste}\NormalTok{(}\StringTok{"x = "}\NormalTok{,}\FunctionTok{sprintf}\NormalTok{(}\StringTok{"\%1.1f"}\NormalTok{, xseq), }\StringTok{"; y = "}\NormalTok{, }\FunctionTok{sprintf}\NormalTok{(}\StringTok{"\%1.1f"}\NormalTok{, yseq),}\StringTok{"; .1"}\NormalTok{, }\AttributeTok{sep=}\StringTok{""}\NormalTok{)}
\FunctionTok{write.table}\NormalTok{(}\FunctionTok{data.frame}\NormalTok{(}\AttributeTok{xyid =}\NormalTok{ xyid, }\AttributeTok{xseq =}\NormalTok{ xseq, }\AttributeTok{yseq =}\NormalTok{ yseq), }\AttributeTok{file =} \FunctionTok{file.path}\NormalTok{(dirODK, }\StringTok{"gridid.csv"}\NormalTok{),}\AttributeTok{row.names =} \ConstantTok{FALSE}\NormalTok{, }\AttributeTok{sep =} \StringTok{","}\NormalTok{)}

\NormalTok{xyid }\OtherTok{\textless{}{-}} \FunctionTok{read.table}\NormalTok{(}\AttributeTok{file =} \FunctionTok{file.path}\NormalTok{(dirODK, }\StringTok{"gridid.csv"}\NormalTok{), }\AttributeTok{header =} \ConstantTok{TRUE}\NormalTok{, }\AttributeTok{sep =} \StringTok{","}\NormalTok{)}
\FunctionTok{str}\NormalTok{(xyid)}
\end{Highlighting}
\end{Shaded}

\subsection{Seleção de subquadrat}\label{seleuxe7uxe3o-de-subquadrat}

O mapa que gera a figura com os subquadrats para serem selecionados.
Preciso ajustar os argumentos para ficar no padrão das outras funções.

\begin{Shaded}
\begin{Highlighting}[]
\FunctionTok{selSubq}\NormalTok{(}\AttributeTok{xmax =} \DecValTok{20}\NormalTok{, }\AttributeTok{ymax =} \DecValTok{20}\NormalTok{, }\AttributeTok{subqX =} \DecValTok{5}\NormalTok{, }\AttributeTok{subqY =} \DecValTok{5}\NormalTok{, }\AttributeTok{mapSize =} \FunctionTok{c}\NormalTok{(}\DecValTok{10}\NormalTok{,}\DecValTok{10}\NormalTok{), }\AttributeTok{fontSize =} \DecValTok{14}\NormalTok{, }\AttributeTok{wd2save =}\NormalTok{ dirODK)}
\end{Highlighting}
\end{Shaded}

\subsection{Mapas e Grids para seleção e mapeamente de
árvores}\label{mapas-e-grids-para-seleuxe7uxe3o-e-mapeamente-de-uxe1rvores}

Para produzir os mapas com as localizações das árvores, precisa criar o
objeto de dados com a função \texttt{mapData}. Aqui estou criando os
dados apenas para a coluna \texttt{A} como exemplo. Os dados de mapa,
assim como os mapas e grids são arquivos grandes que não devem ser
carregados todos de uma vez no dispositivo móvel. Para os grids de
mapeamente, esse objeto não é necessário, já que não há buffer para o
mapeamento.

\begin{Shaded}
\begin{Highlighting}[]
\NormalTok{subplotA }\OtherTok{\textless{}{-}} \FunctionTok{sort}\NormalTok{(}\FunctionTok{grep}\NormalTok{(}\StringTok{"A"}\NormalTok{, }\FunctionTok{unique}\NormalTok{(odkdata}\SpecialCharTok{$}\NormalTok{quad), }\AttributeTok{value =} \ConstantTok{TRUE}\NormalTok{))}
\NormalTok{mapdataA }\OtherTok{\textless{}{-}} \FunctionTok{mapData}\NormalTok{(odkdata[odkdata}\SpecialCharTok{$}\NormalTok{status }\SpecialCharTok{==} \StringTok{"A"}\NormalTok{, ], }\AttributeTok{subplotCodes =}\NormalTok{ subplotA, }\AttributeTok{splitX =} \DecValTok{5}\NormalTok{, }\AttributeTok{splitY =} \DecValTok{5}\NormalTok{, }\AttributeTok{maxX =} \DecValTok{20}\NormalTok{, }\AttributeTok{maxY =} \DecValTok{20}\NormalTok{, }\AttributeTok{buffer =} \DecValTok{2}\NormalTok{)}
\FunctionTok{str}\NormalTok{(mapdataA)}
\NormalTok{nsuqA }\OtherTok{\textless{}{-}} \FunctionTok{length}\NormalTok{(mapdataA)}
\CommentTok{\#subplotA \textless{}{-} paste("A", sprintf("\%02d", 0:15), sep = "")}

\ControlFlowTok{for}\NormalTok{(i }\ControlFlowTok{in}\NormalTok{ subplotA)}
\NormalTok{\{}
    \FunctionTok{svgMap}\NormalTok{(mapdataA, }\AttributeTok{subPlotCode =}\NormalTok{ i, }\AttributeTok{mapSize =} \FunctionTok{c}\NormalTok{(}\DecValTok{12}\NormalTok{, }\DecValTok{12}\NormalTok{), }\AttributeTok{fontSize =} \DecValTok{12}\NormalTok{, }\AttributeTok{wd2save =} \FunctionTok{file.path}\NormalTok{(dirODK, }\StringTok{"subplotA"}\NormalTok{))}
    \FunctionTok{svgGrid}\NormalTok{(}\AttributeTok{censoData =}\NormalTok{ odkdata, }\AttributeTok{subPlotCode =}\NormalTok{ i, }\AttributeTok{mapSize =} \FunctionTok{c}\NormalTok{(}\DecValTok{12}\NormalTok{, }\DecValTok{12}\NormalTok{), }\AttributeTok{fontSize =} \DecValTok{12}\NormalTok{, }\AttributeTok{wd2save =} \FunctionTok{file.path}\NormalTok{(dirODK, }\StringTok{"subplotA"}\NormalTok{))}
    
\NormalTok{\}}
\end{Highlighting}
\end{Shaded}

\section{ODK Collect}\label{odk-collect}

\section{Dispositivo android}\label{dispositivo-android}

O dispositivo android precisa ter a ``Depuração USB'' ativada, para isso
é preciso desbloquear as ``Opções de desenvolvedor'' Conecte o
dispositivo android e garanta que ele tem ``Opções de desenvolvedor'' e
``Depuração USB'' ativados. Caso não esteja, veja como
\href{ativar\%20opção\%20de\%20desenvolvedor}{https://developer.android.com/studio/debug/dev-options?hl=pt-br}.
Entre na ``opções de desenvolvedor'' nas configurações do android e
ative ``Depuração USB''.

\section{Android Debug Bridge}\label{android-debug-bridge}

O ADB é necessário para que o computador conectado ao dispositivo
conecte e transfira arquivos pelo pacote.

No linux (Debian ou Ubuntu) use:

\begin{Shaded}
\begin{Highlighting}[]
\FunctionTok{sudo}\NormalTok{ apt update}
\FunctionTok{sudo}\NormalTok{ apt install android{-}tools{-}adb}
\end{Highlighting}
\end{Shaded}

No archlinux:

\begin{Shaded}
\begin{Highlighting}[]
\FunctionTok{sudo}\NormalTok{ pacman }\AttributeTok{{-}S}\NormalTok{ android{-}tools}
\end{Highlighting}
\end{Shaded}

No windows, consulte:
\href{adb}{https://developer.android.com/tools/adb?hl=pt-br}

\section{Instale e crie o projeto no ODK
Collect}\label{instale-e-crie-o-projeto-no-odk-collect}

Caso a instalação do ODK Collect seja nova, abra o aplicativo e crie um
projeto. A URL para o projeto deve ser: https://odk.ib.usp.br

\section{Conectando o Dispositivo
Android}\label{conectando-o-dispositivo-android}

Após conectar o tablet ao computador abra as notificações e nas opções
de conexão USB selecione a ``Transferir arquivos/Android Auto''.
Mantenha o dispositivo ativo e sem bloqueio de tela.

\subsection{Configurando o USB
Padrão}\label{configurando-o-usb-padruxe3o}

Para tornar o dispositivo android sempre configurado para a troca de
arquivos pelo usb, conecte ao computador, vá para: * Configurações *
Opcções do desenvolvedor * Configurações USB padrão * selecione:
``Transferindo arquivos/Android Auto'' * desconecte o cabo USB e
reconecte * confira a Configuração USB padrão * caso mude para ``Apenas
carregar o tablet'' * selecione novamente ``Tranferindo arquivos/Android
Auto'' * desconecte e reconecte o cabo USB para verificar se a
configuração permanece.

\subsection{Criar projeto no ODK
Collect}\label{criar-projeto-no-odk-collect}

O dispositivo móvel deve ter instalado o app \texttt{ODK\ Collect} e ter
um projeto iniciado. A url do projeto pode ser
\texttt{https:\textbackslash{}\textbackslash{}odk.ib.usp.br}, onde há um
\texttt{ODK\ Central} instalado.

\subsection{Atualizar o formulário
xml}\label{atualizar-o-formuluxe1rio-xml}

Aqui só é necessário atualizar se houve alguma alteração no formulário
\texttt{treeCensusForm.xlsx} depois da instalação do pacote.

O formulário é criado em xlsx no libreOffice. É preciso traduzir para
formato xml utilizando o pacote do python pyxform. O pyxform foi
incorporado no pacote do R, assim como o \texttt{treeCensusForm.xlsx} e
o \texttt{treeCensusForm.xml}.

\begin{tcolorbox}[enhanced jigsaw, arc=.35mm, bottomrule=.15mm, bottomtitle=1mm, breakable, colback=white, colbacktitle=quarto-callout-warning-color!10!white, colframe=quarto-callout-warning-color-frame, coltitle=black, left=2mm, leftrule=.75mm, opacityback=0, opacitybacktitle=0.6, rightrule=.15mm, title=\textcolor{quarto-callout-warning-color}{\faExclamationTriangle}\hspace{0.5em}{Atenção}, titlerule=0mm, toprule=.15mm, toptitle=1mm]

Este procedimento é dispensável, o pacote já é distribuído com o
formulário xml traduzido e atualizado. Ele só deve ser usado se houver
alguma mudança local no xlsxform (planilha LibreOffice) que precisa ser
incorporada no formulário, em casos de mudanças locais emergenciais.
Caso seja um bug, deve ser informado no isssues do github para a solução
seja incorporada na distribuição do pacote.

\end{tcolorbox}

\begin{Shaded}
\begin{Highlighting}[]
\NormalTok{xlsxFile }\OtherTok{\textless{}{-}} \FunctionTok{system.file}\NormalTok{(}\StringTok{"odk/treeCensusForm.xlsx"}\NormalTok{, }\AttributeTok{package =} \StringTok{"odkTreeCensus"}\NormalTok{)}
\NormalTok{pkgDirForm }\OtherTok{\textless{}{-}} \FunctionTok{dirname}\NormalTok{(xlsxFile)}
\NormalTok{xmlFile }\OtherTok{\textless{}{-}} \FunctionTok{file.path}\NormalTok{(pkgDirForm, }\StringTok{"treeCensusForm.xml"}\NormalTok{)}
\FunctionTok{xmlForm}\NormalTok{(}\AttributeTok{xlsxPath =}\NormalTok{ xlsxFile, }\AttributeTok{xmlPath =}\NormalTok{ xmlFile)}
\end{Highlighting}
\end{Shaded}

\section{Copia o caminho do formulário
XML}\label{copia-o-caminho-do-formuluxe1rio-xml}

Copiando o caminho onde se encontra o formulário \texttt{xml} do pacote
para tranferir ao dispositivo móvel.

\begin{Shaded}
\begin{Highlighting}[]
\NormalTok{xmlDir }\OtherTok{\textless{}{-}} \FunctionTok{system.file}\NormalTok{(}\StringTok{"odk"}\NormalTok{, }\AttributeTok{package =} \StringTok{"odkTreeCensus"}\NormalTok{)}
\NormalTok{xmlPath }\OtherTok{\textless{}{-}} \FunctionTok{system.file}\NormalTok{(}\StringTok{"odk/treeCensusForm.xml"}\NormalTok{, }\AttributeTok{package =} \StringTok{"odkTreeCensus"}\NormalTok{)}
\end{Highlighting}
\end{Shaded}

\section{Subindo o formulário para o dispositivo móvel
android}\label{subindo-o-formuluxe1rio-para-o-dispositivo-muxf3vel-android}

\begin{Shaded}
\begin{Highlighting}[]
\FunctionTok{pushToODKCollect}\NormalTok{(}\AttributeTok{localDir =}\NormalTok{ xmlDir, }\AttributeTok{formName =} \StringTok{"treeCensusForm"}\NormalTok{)}
\FunctionTok{mediaToODK}\NormalTok{(}\AttributeTok{localMediaDir =}\NormalTok{ dirODK, }\AttributeTok{formName =} \StringTok{"treeCensusForm"}\NormalTok{)}
\end{Highlighting}
\end{Shaded}





\end{document}
